\documentclass[11pt]{article}

\usepackage[T1]{fontenc}
\usepackage[utf8]{inputenc}
\usepackage[english]{babel}
\usepackage[margin=1in]{geometry}
\usepackage{lmodern}
\usepackage{microtype}
\usepackage{mathtools}
\usepackage{amsmath,amssymb,amsthm}
\usepackage{array}
\usepackage{booktabs}
\usepackage{tabularx}
\usepackage{longtable}
\usepackage{float}
\usepackage{xcolor}
\usepackage{graphicx}
\graphicspath{{contents/figures/}{figures/}}
\usepackage{pdflscape}
\usepackage{caption}
\usepackage{enumitem}
\usepackage{fancyhdr}
\usepackage[numbers]{natbib}
\usepackage{doi}
\usepackage{xurl}
\usepackage{hyperref}
\usepackage{cleveref}
\usepackage{orcidlink}
\usepackage{titling}

\input{includes/paper_macros.tex}
\providecommand{\tightlist}{%
  \setlength{\itemsep}{0pt}\setlength{\parskip}{0pt}%
}
\newif\ifpaperhascitations
\paperhascitationstrue

\newtheorem{theorem}{Theorem}[section]
\newtheorem{lemma}[theorem]{Lemma}
\newtheorem{proposition}[theorem]{Proposition}
\newtheorem{corollary}[theorem]{Corollary}
\theoremstyle{definition}
\newtheorem{definition}[theorem]{Definition}
\newtheorem{remark}[theorem]{Remark}
\floatstyle{ruled}
\newfloat{algorithm}{tbp}{loa}
\floatname{algorithm}{Algorithm}

\urlstyle{same}
\captionsetup{
  font=small,
  labelfont=bf,
  labelsep=period
}
\captionsetup[table]{skip=6pt,position=top}
\captionsetup[figure]{skip=6pt}
\captionsetup[algorithm]{skip=6pt,position=top}
\crefname{algorithm}{algorithm}{algorithms}
\Crefname{algorithm}{Algorithm}{Algorithms}
\crefname{theorem}{theorem}{theorems}
\Crefname{theorem}{Theorem}{Theorems}
\crefname{lemma}{lemma}{lemmas}
\Crefname{lemma}{Lemma}{Lemmas}
\crefname{proposition}{proposition}{propositions}
\Crefname{proposition}{Proposition}{Propositions}
\crefname{corollary}{corollary}{corollaries}
\Crefname{corollary}{Corollary}{Corollaries}
\crefname{definition}{definition}{definitions}
\Crefname{definition}{Definition}{Definitions}
\crefname{remark}{remark}{remarks}
\Crefname{remark}{Remark}{Remarks}
\setlist[enumerate]{leftmargin=*, itemsep=2pt, topsep=4pt}
\setlength{\droptitle}{-3.2em}
\pretitle{\begin{center}\LARGE\bfseries}
\posttitle{\par\end{center}\vspace{0.5em}}
\preauthor{\begin{center}\normalsize}
\postauthor{\par\end{center}\vspace{-0.4em}}
\predate{\begin{center}\small}
\postdate{\par\end{center}\vspace{0.6em}}
\setlength{\emergencystretch}{2em}
\clubpenalty=10000
\widowpenalty=10000
\displaywidowpenalty=10000
\interfootnotelinepenalty=10000
\allowdisplaybreaks[2]
\fancypagestyle{firstpage}{\fancyhf{}
  \fancyfoot[C]{\scriptsize Automorph Inc., Wilmington, DE, USA \quad \texttt{ioannis@automorph.io} \quad \textcopyright\ 2026 \quad CC-BY 4.0 \quad \doi{10.5281/zenodo.23218104}}
  \renewcommand{\headrulewidth}{0pt}
  \renewcommand{\footrulewidth}{0pt}
}

\title{The Cosmological Constant as a Contingent Modulus:\\ A Six Birds Demarcation of Dark Energy}
\author{Ioannis Tsiokos\,\orcidlink{0009-0009-7659-5964}}
\date{Version 1: 17 June 2026\quad$\cdot$\quad Version 2: 7 October 2026}
\hypersetup{
  pdftitle={The Cosmological Constant as a Contingent Modulus: A Six Birds Demarcation of Dark Energy},
  pdfauthor={Ioannis Tsiokos},
  hidelinks
}

\begin{document}

\maketitle
\thispagestyle{firstpage}

\begin{abstract}
Dark energy poses several different questions at once: what kind of term
the cosmological constant is, what fixes its value, why its density is
comparable to that of matter today, and what current surveys say about
an evolving equation of state. We study these questions in
finite-dimensional mathematical models built with the Six Birds closure
calculus, and for each model we separate what it determines from what it
leaves as input. First, in a convex closure problem that attains its
optimum and whose optimal value is finite and differentiable near the
budget level, a budget whose relaxation strictly lowers the optimal cost
carries a single positive price; in an allocation model with an
originally binding budget, a common linear utility term leaves the
optimizer unchanged and shifts only that price, as long as the shifted
price stays nonnegative. Second, in a self-consistent lattice model, the
convergence readouts used to admit background values do not determine
the background: all eighteen sampled offsets converge, fifteen pass the
sixty-update test, and neither readout recovers the input. When a global
density ledger is added, the absolute background becomes a gauge
coordinate that the invariant record cannot recover, and no
translation-invariant probability distribution over it exists. Third, in
a toy universe with structure growth and record erasure, at least ten of
twelve sampled populations retain distinguishable records exactly for
\(-1/8000<\lambda\le\lambda_+\), with \(3<\lambda_+<3.4\), but the ratio
\(|\lambda|/\rho_m\) at the first record is not pinned near unity.
Whether the present epoch is typical then depends on a declared prior:
none of five tested priors meets a declared concentration criterion, and
smoothness near zero is not enough. On DESI DR2, Pantheon+ and
compressed Planck data, pre-registered audits find a fitted rewrite
shape compatible with a constant, show that the best-fit evolving
history has a positive, regular density through its \(w=-1\) crossing,
and flag structured dependence between supernova and BAO subsets, led by
high-redshift supernovae. The statements labelled as Lean theorems are
kernel-checked. None of these results identifies a model price with the
physical cosmological constant or derives its value.
\end{abstract}

\medskip
\noindent\textbf{Keywords.} dark energy; cosmological constant; shadow prices; identifiability; coincidence problem; Lean; emergence calculus; Six Birds Theory.

\input{sections/sec_01_introduction}
\input{sections/sec_02_the_calculus}
\input{sections/sec_03_method}
\input{sections/sec_04_form}
\input{sections/sec_05_window}
\input{sections/sec_06_value}
\input{sections/sec_07_measure}
\input{sections/sec_08_audits}
\input{sections/sec_09_mysteries}
\input{sections/sec_10_related}
\input{sections/sec_11_scope}
\input{sections/sec_12_conclusion}

\appendix
\input{appendices/app_a_carriers}
\input{appendices/app_b_audits}
\input{appendices/app_c_reproducibility}
\input{appendices/app_d_review}
\input{appendices/app_e_version_history}

\ifpaperhascitations
\clearpage
\bibliographystyle{plainnat}
\bibliography{references}
\fi

\end{document}
